\documentclass[12pt,a4paper]{article}

% --- Packages ---
\usepackage[margin=2.5cm]{geometry}
\usepackage{amsmath,amssymb,amsthm}
\usepackage{mathtools}
\usepackage{booktabs}
\usepackage{array}
\usepackage{xcolor}
\usepackage{hyperref}
\usepackage{enumitem}
\usepackage{titlesec}
\usepackage{fancyhdr}
\usepackage{graphicx}
\usepackage{tcolorbox}
\usepackage{parskip}
\usepackage{microtype}
\usepackage{float}

% --- Theorem environments ---
\newtheorem{theorem}{Theorem}
\newtheorem{lemma}[theorem]{Lemma}
\newtheorem{proposition}[theorem]{Proposition}
\newtheorem{definition}{Definition}
\newtheorem{remark}{Remark}

% --- Colors ---
\definecolor{deepblue}{RGB}{0,56,101}
\definecolor{lightblue}{RGB}{220,235,250}
\definecolor{boxgray}{RGB}{245,245,245}

% --- Section formatting ---
\titleformat{\section}{\large\bfseries\color{deepblue}}{}{0em}{\thesection\quad}
\titleformat{\subsection}{\normalsize\bfseries\color{deepblue}}{}{0em}{\thesubsection\quad}

% --- Header/footer ---
\pagestyle{fancy}
\fancyhf{}
\rhead{\small Drone-Based Delivery Optimization}
\lhead{\small ENSIA -- Third Year -- 2026}
\cfoot{\thepage}
\renewcommand{\headrulewidth}{0.4pt}

% --- Hyperref setup ---
\hypersetup{
  colorlinks=true,
  linkcolor=deepblue,
  urlcolor=deepblue,
  citecolor=deepblue
}

% --- Header height fix ---
\setlength{\headheight}{15pt}

% --- Shorthand math commands ---
\newcommand{\R}{\mathbb{R}}
\newcommand{\Z}{\mathbb{Z}}
\newcommand{\Nset}{\mathcal{N}}
\newcommand{\Kset}{\mathcal{K}}

% =====================================================================
\begin{document}
% =====================================================================

% --- Title page ---
\begin{titlepage}
  \centering
  \vspace*{2cm}
  {\huge\bfseries\color{deepblue} Drone-Based Delivery Optimization}\\[0.6em]
  {\large\bfseries Project Report}\\[2em]
  {\large National School of Artificial Intelligence (ENSIA)}\\[0.4em]
  {\large Third Year -- Combinatorial Optimization and Metaheuristics}\\[0.4em]
  {\large Academic Year 2025--2026}\\[4em]
  \begin{tcolorbox}[colback=lightblue,colframe=deepblue,width=0.85\textwidth]
    \centering
    \textbf{Project 3:} Drone-Based Delivery Optimization Problem\\[0.4em]
    Submission deadline: Saturday, 16 May 2026
  \end{tcolorbox}
  \vspace{2em}
  \begin{tcolorbox}[colback=boxgray,colframe=deepblue,width=0.6\textwidth]
    \centering
    \textbf{Group Members}\\[0.5em]
    BRIHMOUCHE Aicha\\
    LALAM Sofia
  \end{tcolorbox}
  \vfill
  {\small This report presents: two distinct mathematical formulations, a complexity study,
  an exact Branch-and-Bound method, two metaheuristics (Genetic Algorithm and Simulated
  Annealing), and a comparative experimental study on multiple benchmark instances.}
\end{titlepage}

\tableofcontents
\newpage

% =====================================================================
\section{Problem Description and Understanding}
% =====================================================================

\subsection{Context and Motivation}

We consider a drone-based last-mile delivery system in which a logistics company operates
a fleet of identical drones to deliver packages from a central depot to a set of customers.
This problem is motivated by the growing interest in autonomous aerial delivery in
resource-constrained settings, such as Algeria, where drone fleets are small and battery
technology limits flight range.

Because the number of drones is limited and battery capacity restricts each drone to a
finite flight range, it is generally \emph{impossible} to serve every customer in a single
planning horizon. Rather than attempting universal service and failing, we adopt the
philosophy of the \textbf{Team Orienteering Problem (TOP)}: prioritize the most urgent
customers and design routes that maximize the total collected urgency/profit while
respecting physical constraints.

\subsection{Problem Description}

The problem can be informally stated as follows:

\begin{tcolorbox}[colback=boxgray,colframe=deepblue]
Given a depot, a set of customers each with a demand and an urgency score, and a fleet of
$m$ identical drones, determine which customers to serve and in which order, such that:
\begin{itemize}[leftmargin=*]
  \item each drone starts and ends its route at the depot;
  \item each drone's total energy consumption does not exceed its battery capacity $B$;
  \item each drone's total carried payload does not exceed its weight capacity $Q$;
  \item each customer is visited at most once across all drones;
  \item the total collected urgency/profit is maximized.
\end{itemize}
\end{tcolorbox}

\subsection{Problem Variants Combined}

Our formulation combines elements from four classical routing problems:

\begin{itemize}
  \item \textbf{Vehicle Routing Problem (VRP):} multi-route planning with a depot.
  \item \textbf{Capacitated VRP (CVRP):} payload capacity $Q$ per drone per route.
  \item \textbf{Electric VRP (EVRP):} energy/battery constraints replacing fuel.
  \item \textbf{Team Orienteering Problem (TOP):} selective customer visits to maximize
  profit under a travel-budget constraint.
\end{itemize}

\subsection{Problem Assumptions}

\begin{enumerate}
  \item The drone fleet is \textbf{homogeneous}: all drones are identical.
  \item All parameters are \textbf{deterministic}: demands, distances, and energy
  consumption are known and fixed.
  \item Loading and unloading at any node takes \textbf{negligible time}.
  \item Each drone performs exactly \textbf{one route} per battery charge; it may return
  to the depot for reloading if battery permits, but no mid-route recharging occurs.
  \item Distances are \textbf{Euclidean} (or can be replaced by GIS-based distances).
  \item There are no airspace conflicts between simultaneous drone operations.
  \item Energy consumption on arc $(i,j)$ depends linearly on distance and current
  payload (consistent with multirotor physics):
  \[
    e_{ij}(\ell) = (\alpha \cdot (m_d + m_b + \ell) + \beta) \cdot d_{ij}
  \]
  where $m_d$ is the drone mass, $m_b$ is the battery mass, $\ell$ is the current
  payload, and $\alpha, \beta > 0$ are drone-specific constants.
\end{enumerate}

\subsection{Customer Prioritization: The TOP Pre-Processing Step}

Since resources are scarce, we introduce a pre-processing step that selects a tractable
subset of high-priority customers before solving the routing problem.

\begin{definition}[Priority Subset Selection]
Let $S = \{1, 2, \ldots, s\}$ be the full set of potential customers. Each customer
$i \in S$ is assigned an urgency score $p_i \geq 0$, where higher values indicate
higher delivery priority (e.g., medical supplies $>$ food $>$ books). After sorting
customers in decreasing order of $p_i$, we select the top $K_{\max}$ customers that
can realistically be served given fleet capacity. Let $N \subseteq S$ denote this
selected subset.
\end{definition}

After this selection, the optimization model is solved only on the reduced set $N$.

% =====================================================================
\section{Mathematical Notation}
% =====================================================================

The following sets, parameters, and decision variables are shared by both formulations.

\subsection{Sets}

\begin{center}
\begin{tabular}{ll}
\toprule
Symbol & Description \\
\midrule
$N = \{1, 2, \ldots, n\}$ & Set of selected customers \\
$N_0 = N \cup \{0\}$ & Customers plus depot (node 0) \\
$K = \{1, 2, \ldots, m\}$ & Set of identical drones \\
\bottomrule
\end{tabular}
\end{center}

\subsection{Parameters}

\begin{center}
\begin{tabular}{ll}
\toprule
Symbol & Description \\
\midrule
$d_{ij}$ & Euclidean distance from node $i$ to node $j$ \\
$e_{ij}$ & Energy consumed flying arc $(i,j)$ at average payload \\
$q_i$ & Demand (payload weight) of customer $i$; $q_0 = 0$ \\
$p_i$ & Urgency/profit score of customer $i$ \\
$B$ & Battery capacity per drone (energy units) \\
$Q$ & Payload capacity per drone (weight units) \\
$m^{\max}_k$ & Maximum number of trips allowed for drone $k$ \\
$\lambda$ & Small energy penalty weight ($0.01 \leq \lambda \leq 0.1$) \\
$M$ & A sufficiently large constant (Big-M) \\
\bottomrule
\end{tabular}
\end{center}

% =====================================================================
\section{Formulation 1: Arc-Based Mixed Integer Linear Program}
% =====================================================================

\subsection{Modeling Philosophy}

The first formulation is a \textbf{three-index arc-based MILP}. The central decision is
whether drone $k$ traverses the arc from node $i$ to node $j$. This is the classical
approach for vehicle routing problems, making battery and payload tracking explicit
through auxiliary continuous variables on each arc.

\subsection{Decision Variables}

\begin{center}
\begin{tabular}{ll}
\toprule
Variable & Description \\
\midrule
$x_{ijk} \in \{0,1\}$ & 1 if drone $k$ travels directly from $i$ to $j$ \\
$y_i \in \{0,1\}$ & 1 if customer $i$ is served by any drone \\
$u_{ik} \geq 0$ & Remaining payload on drone $k$ upon departure from node $i$ \\
$b_{ik} \geq 0$ & Remaining battery of drone $k$ upon arrival at node $i$ \\
$z_k \in \{0,1\}$ & 1 if drone $k$ is activated (used) \\
\bottomrule
\end{tabular}
\end{center}

\subsection{Objective Function}

We maximize a weighted combination of total urgency collected minus a small penalty on
total energy:

\begin{equation}
  \max \quad \sum_{i \in N} p_i \, y_i \;-\; \lambda \sum_{k \in K} \sum_{i \in N_0}
  \sum_{\substack{j \in N_0 \\ j \neq i}} e_{ij} \, x_{ijk}
  \tag{F1-Obj}
\end{equation}

The first term maximizes total urgency/profit; the second penalizes energy use with
weight $\lambda \ll 1$ so that profit remains the dominant objective.

\subsection{Constraints}

\paragraph{(C1) Customer serving -- link $y_i$ to arc variables:}
\begin{equation}
  \sum_{k \in K} \sum_{\substack{j \in N_0 \\ j \neq i}} x_{ijk} = y_i
  \qquad \forall\, i \in N
  \tag{C1}
\end{equation}
Customer $i$ is served (arrived at) exactly as many times as $y_i$ indicates. Combined
with flow conservation, this guarantees each customer is visited at most once.

\paragraph{(C2) Flow conservation per drone at each node:}
\begin{equation}
  \sum_{\substack{i \in N_0 \\ i \neq j}} x_{ijk} =
  \sum_{\substack{l \in N_0 \\ l \neq j}} x_{jlk}
  \qquad \forall\, j \in N_0,\; k \in K
  \tag{C2}
\end{equation}
Every time drone $k$ enters node $j$, it must also leave. This holds at the depot too,
allowing multiple trips (the drone returns and departs again).

\paragraph{(C3) Maximum trips per drone:}
\begin{equation}
  \sum_{j \in N} x_{0jk} \leq m^{\max}_k
  \qquad \forall\, k \in K
  \tag{C3}
\end{equation}
Limits the number of times drone $k$ departs from the depot.

\paragraph{(C4) Initial battery -- full at first departure:}
\begin{equation}
  b_{0k} = B
  \qquad \forall\, k \in K
  \tag{C4}
\end{equation}

\paragraph{(C5) Battery consumption tracking (Big-M formulation):}
\begin{equation}
  b_{ik} - e_{ij} \geq b_{jk} - M(1 - x_{ijk})
  \qquad \forall\, i,j \in N_0,\; i \neq j,\; k \in K
  \tag{C5}
\end{equation}
If drone $k$ travels arc $(i,j)$ (i.e., $x_{ijk} = 1$), then the battery at $j$ equals
the battery at $i$ minus the arc's energy cost. When $x_{ijk} = 0$, the Big-M makes
the constraint non-binding.

\paragraph{(C6) Battery bounds:}
\begin{equation}
  0 \leq b_{ik} \leq B
  \qquad \forall\, i \in N_0,\; k \in K
  \tag{C6}
\end{equation}

\paragraph{(C7) Initial payload -- full at depot:}
\begin{equation}
  u_{0k} = Q
  \qquad \forall\, k \in K
  \tag{C7}
\end{equation}

\paragraph{(C8) Payload consumption tracking:}
\begin{equation}
  u_{ik} - q_j \geq u_{jk} - M(1 - x_{ijk})
  \qquad \forall\, i \in N_0,\; j \in N,\; i \neq j,\; k \in K
  \tag{C8}
\end{equation}
Payload at $j$ equals payload at $i$ minus the demand delivered at $j$, when
$x_{ijk} = 1$.

\paragraph{(C9) Payload bounds:}
\begin{equation}
  0 \leq u_{ik} \leq Q
  \qquad \forall\, i \in N_0,\; k \in K
  \tag{C9}
\end{equation}

\paragraph{(C10) At least one customer served:}
\begin{equation}
  \sum_{i \in N} y_i \geq 1
  \tag{C10}
\end{equation}

\paragraph{(C11) Variable domains:}
\begin{equation}
  x_{ijk} \in \{0,1\}, \quad y_i \in \{0,1\}, \quad z_k \in \{0,1\}
  \tag{C11}
\end{equation}

\subsection{Formulation Summary}

The complete model (F1) is:
\begin{align*}
  \max \quad & \text{(F1-Obj)}  \\
  \text{s.t.} \quad & \text{(C1)--(C11)}
\end{align*}

\noindent\textbf{Model size.} With $n$ customers and $m$ drones:
\begin{itemize}
  \item Binary variables: $O(n^2 m)$ for $x_{ijk}$, $O(n)$ for $y_i$, $O(m)$ for $z_k$.
  \item Continuous variables: $O(nm)$ for $u_{ik}$ and $b_{ik}$.
  \item Constraints: $O(n^2 m)$ mainly from (C5) and (C8).
\end{itemize}

% =====================================================================
\section{Formulation 2: Route-Based Set Partitioning Model}
% =====================================================================

\subsection{Modeling Philosophy}

The second formulation takes a fundamentally different view of the problem. Instead of
tracking arc-by-arc movements, we model the problem at the level of \textbf{complete
drone routes}. A route is a full round trip: depot $\to$ some customers $\to$ depot.
We enumerate (or generate) all routes that are \emph{feasible by construction} -- i.e.,
routes that already respect battery and payload limits -- and then ask which subset of
routes to select.

This transforms the routing problem into a \textbf{set partitioning / set covering
integer program}, a classical paradigm in vehicle routing literature [2,3].

The key contrast with Formulation 1:

\begin{center}
\renewcommand{\arraystretch}{1.4}
\begin{tabular}{lll}
\toprule
Aspect & Formulation 1 (Arc-Based) & Formulation 2 (Route-Based) \\
\midrule
Decision unit & Individual arc $(i \to j)$ by drone $k$ & Complete route $r$ \\
Feasibility & Enforced by constraints (C5)--(C9) & Baked into route definition \\
Subtour elim. & Needs explicit constraints & Not needed \\
Variables & $O(n^2 m)$ binary + continuous & One binary per feasible route \\
Battery tracking & Explicit $b_{ik}$ variables & Implicit in route energy check \\
LP bound quality & Moderate & Tighter (fewer fractional solutions) \\
\bottomrule
\end{tabular}
\end{center}

\subsection{The Route Set $\Omega$}

\begin{definition}[Feasible Route]
A \textbf{feasible route} $r$ is a sequence of nodes
$r = (0, i_1, i_2, \ldots, i_\ell, 0)$ where:
\begin{enumerate}
  \item The route starts and ends at the depot (node 0).
  \item Each customer node appears at most once: $\{i_1, \ldots, i_\ell\} \subseteq N$,
  all distinct.
  \item The total energy consumed does not exceed battery capacity:
  \[
    E(r) = \sum_{t=0}^{\ell} e_{i_t, i_{t+1}} \leq B
    \qquad (i_0 = i_{\ell+1} = 0)
  \]
  \item The total payload carried at departure does not exceed drone capacity:
  \[
    W(r) = \sum_{t=1}^{\ell} q_{i_t} \leq Q
  \]
\end{enumerate}
Let $\Omega$ denote the set of all feasible routes.
\end{definition}

\noindent For each route $r \in \Omega$, define the following route parameters:
\begin{itemize}
  \item $P(r) = \sum_{i \in r} p_i$ \quad -- total profit/urgency collected on route $r$.
  \item $E(r)$ \quad -- total energy consumed on route $r$ (as defined above).
  \item $a_{ir} \in \{0,1\}$ \quad -- 1 if customer $i$ is visited on route $r$, 0
  otherwise.
\end{itemize}

\subsection{Decision Variable}

\begin{center}
\begin{tabular}{ll}
\toprule
Variable & Description \\
\midrule
$\lambda_r \in \{0,1\}$ & 1 if route $r \in \Omega$ is selected and assigned to a drone \\
\bottomrule
\end{tabular}
\end{center}

\subsection{Objective Function}

\begin{equation}
  \max \quad \sum_{r \in \Omega} \left( P(r) - \delta \cdot E(r) \right) \lambda_r
  \tag{F2-Obj}
\end{equation}

where $\delta > 0$ is a small penalty on energy, analogous to $\lambda$ in Formulation 1.
The objective selects a set of routes that collectively maximizes profit while penalizing
energy-intensive routes.

\subsection{Constraints}

\paragraph{(D1) Each customer served at most once:}
\begin{equation}
  \sum_{r \in \Omega} a_{ir} \, \lambda_r \leq 1
  \qquad \forall\, i \in N
  \tag{D1}
\end{equation}
No customer appears in more than one selected route. This is the \emph{set covering}
constraint that gives this formulation its name.

\paragraph{(D2) Fleet size -- at most $m$ routes selected:}
\begin{equation}
  \sum_{r \in \Omega} \lambda_r \leq m
  \tag{D2}
\end{equation}
Since we have exactly $m$ drones and each drone executes at most one route.

\paragraph{(D3) At least one route selected:}
\begin{equation}
  \sum_{r \in \Omega} \lambda_r \geq 1
  \tag{D3}
\end{equation}

\paragraph{(D4) Binary variable domain:}
\begin{equation}
  \lambda_r \in \{0,1\}
  \qquad \forall\, r \in \Omega
  \tag{D4}
\end{equation}

\subsection{Complete Formulation (F2)}

\begin{align*}
  \max \quad & \sum_{r \in \Omega} \left( P(r) - \delta \cdot E(r) \right) \lambda_r \\
  \text{s.t.} \quad
  & \sum_{r \in \Omega} a_{ir} \, \lambda_r \leq 1 \qquad \forall\, i \in N \\
  & \sum_{r \in \Omega} \lambda_r \leq m \\
  & \sum_{r \in \Omega} \lambda_r \geq 1 \\
  & \lambda_r \in \{0,1\} \qquad \forall\, r \in \Omega
\end{align*}

\subsection{Why This is a Genuinely Different Formulation}

Formulation 1 and Formulation 2 represent two fundamentally different mathematical
abstractions of the same physical problem:

\begin{itemize}
  \item \textbf{Abstraction level.} Formulation 1 operates at the \emph{arc} level --
  the finest granularity. Each binary variable records a single drone movement between
  two nodes. Formulation 2 operates at the \emph{route} level -- a higher abstraction
  where each variable encodes an entire drone mission.

  \item \textbf{Constraint structure.} In Formulation 1, battery and payload feasibility
  is enforced by a large number of explicit Big-M constraints (C5) and (C8). In
  Formulation 2, feasibility is embedded in the \emph{definition} of $\Omega$: only
  routes that satisfy battery and payload limits are in $\Omega$. There are no Big-M
  constraints at all.

  \item \textbf{Subtour elimination.} Formulation 1 requires either MTZ constraints or
  subtour elimination cuts to prevent infeasible loops. Formulation 2 requires none:
  every route in $\Omega$ is by definition a simple path starting and ending at the depot.

  \item \textbf{LP relaxation quality.} The LP relaxation of Formulation 2 is known
  to provide tighter bounds than the LP relaxation of Formulation 1, because the route
  variables implicitly capture the combinatorial structure of feasible tours [3]. This is the main theoretical advantage of the route-based model.

  \item \textbf{Computational trade-off.} Formulation 1 has a polynomial number of
  variables but weaker LP bounds; it is directly solvable by a MILP solver.
  Formulation 2 has an exponential number of variables (one per feasible route) but
  stronger LP bounds; it is typically solved via \emph{column generation}, where routes
  (columns) are generated on demand by solving a shortest-path pricing sub-problem.
  For our project's small instance sizes, we can enumerate $\Omega$ explicitly.
\end{itemize}

\subsection{Practical Generation of $\Omega$}

For small instances (up to 10--15 customers), we generate $\Omega$ explicitly:

\begin{enumerate}
  \item Enumerate all subsets $S \subseteq N$ with $\sum_{i \in S} q_i \leq Q$.
  \item For each subset $S$, find the shortest Hamiltonian path visiting all nodes in $S$
  starting and ending at the depot (solve a small TSP).
  \item Check if the energy of the best ordering satisfies $E \leq B$.
  \item If yes, add the route to $\Omega$.
\end{enumerate}

For larger instances, column generation is used: only routes with negative reduced cost
are generated at each iteration of the LP relaxation.

% =====================================================================
\section{Problem Complexity}
\label{sec:complexity}
% =====================================================================

\subsection{Statement}

\begin{theorem}
The Drone-Based Delivery Optimization Problem (DBDO) is NP-hard.
\end{theorem}

\subsection{Membership in NP}

We first show that the decision version of DBDO is in NP.

\begin{definition}[Decision version of DBDO]
Given an instance of DBDO and bounds $P^* \geq 0$ and $E^* \geq 0$: does there exist
a set of drone routes such that the total collected profit is at least $P^*$ and the
total energy consumed is at most $E^*$?
\end{definition}

\begin{lemma}
DBDO (decision version) is in NP.
\end{lemma}

\begin{proof}
A certificate for a ``yes'' instance is a set of at most $m$ routes
$\{r_1, r_2, \ldots, r_m\}$, where each $r_k$ is a sequence of nodes. Given this
certificate, a verifier can:
\begin{enumerate}
  \item Check that each route starts and ends at node 0. \hfill $O(n)$ per route.
  \item Check that no customer appears in more than one route. \hfill $O(nm)$.
  \item Compute the energy of each route and check $E(r_k) \leq B$ for all $k$.
  \hfill $O(n)$ per route.
  \item Compute the payload of each route and check $W(r_k) \leq Q$ for all $k$.
  \hfill $O(n)$ per route.
  \item Sum the profits and check $\sum_k P(r_k) \geq P^*$. \hfill $O(n)$.
\end{enumerate}
The total verification time is $O(n^2)$, which is polynomial. Therefore DBDO is in NP.
\end{proof}

\subsection{NP-Hardness via Reduction from TSP}

We prove NP-hardness by showing that the classical Traveling Salesman Problem (TSP)
reduces in polynomial time to DBDO. Since TSP is NP-hard, this implies DBDO is NP-hard.

\subsubsection{Recall: The TSP Decision Problem}

Given a complete undirected graph $G = (V, E)$ with edge costs $c_{ij} \geq 0$ for
$i,j \in V$ with $|V| = n$, and a bound $C \geq 0$:

\begin{center}
\emph{Does there exist a Hamiltonian cycle of total cost at most $C$?}
\end{center}

TSP (decision version) is NP-complete [4].

\subsubsection{Polynomial Reduction: TSP $\leq_p$ DBDO}

Given an arbitrary TSP instance $(G, c, C)$, we construct a DBDO instance as follows:

\medskip
\noindent\textbf{Construction} (polynomial time, $O(n^2)$):

\begin{itemize}
  \item Set $m = 1$ (one drone only).
  \item Use the same node set: place the depot at an arbitrary node $v_0 \in V$, and set
  customers $N = V \setminus \{v_0\}$, so $|N| = n-1$.
  \item Set energy costs equal to TSP distances: $e_{ij} = c_{ij}$ for all $i,j$.
  \item Set battery capacity $B$ to a large number: $B = \sum_{i,j} c_{ij}$ (effectively
  unlimited; this is a constant computable in $O(n^2)$).
  \item Set payload capacity $Q$ to a large number: $Q = n$ (effectively unlimited).
  \item Set all demands $q_i = 0$ for all $i \in N$ (no payload constraint).
  \item Set all urgency scores $p_i = 1$ for all $i \in N$.
  \item Set energy penalty $\lambda = 1$.
  \item Set threshold $P^* = n-1$ (require all customers to be served).
  \item Set energy bound $E^* = C$.
\end{itemize}

This construction takes $O(n^2)$ time and space.

\subsubsection{Correctness: If and Only If}

\begin{proof}[Proof of Theorem 1]
We prove both directions of the equivalence.

\medskip
\noindent($\Rightarrow$) \textbf{TSP has a tour of cost $\leq C$ $\Rightarrow$ DBDO
has a solution with profit $\geq P^*$ and energy $\leq E^*$.}

Suppose the TSP instance has a Hamiltonian cycle $v_0 \to v_1 \to \cdots \to v_{n-1}
\to v_0$ with total cost $\leq C$. We construct a DBDO solution directly: use this
cycle as the single drone's route. Specifically, the route is the sequence
$0 \to i_1 \to i_2 \to \cdots \to i_{n-1} \to 0$, where $i_t = v_t$.

\begin{itemize}
  \item \emph{Feasibility:} Since $B$ is unlimited, the energy constraint $E(r) \leq B$
  is trivially satisfied. Since $Q$ is unlimited and all $q_i = 0$, the payload
  constraint is trivially satisfied.
  \item \emph{Profit:} The route visits all $n-1$ customers, each with profit 1, so
  total profit $= n-1 = P^*$. \checkmark
  \item \emph{Energy:} The energy equals the TSP tour cost $\leq C = E^*$. \checkmark
\end{itemize}

Therefore the DBDO instance has a feasible solution meeting both bounds.

\medskip
\noindent($\Leftarrow$) \textbf{DBDO has a solution with profit $\geq P^* = n-1$ and
energy $\leq E^* = C$ $\Rightarrow$ TSP has a tour of cost $\leq C$.}

Suppose the DBDO instance has a feasible solution. Since $P^* = n-1$ and each customer
contributes profit exactly 1, the solution must visit \emph{every} customer in $N$
(visiting fewer gives profit $< n-1 = P^*$, which violates the profit requirement).

Since battery and payload are unlimited ($B$ and $Q$ are very large) and there is only
one drone, the solution is a single route that:
\begin{enumerate}
  \item starts at the depot $v_0$,
  \item visits all $n-1$ customers exactly once (as shown above),
  \item returns to the depot.
\end{enumerate}

This is precisely a Hamiltonian cycle of the original graph $G$. The energy of this
route equals $\sum_{(i,j) \in \text{route}} c_{ij} \leq E^* = C$.

Therefore the TSP instance has a Hamiltonian cycle of cost $\leq C$. \checkmark

\medskip
We have shown: TSP has a tour of cost $\leq C$ \emph{if and only if} the constructed
DBDO instance has a solution with profit $\geq n-1$ and energy $\leq C$. Since TSP is
NP-hard and this reduction is polynomial, DBDO is NP-hard. Combined with the NP
membership result from Lemma 1, DBDO is \textbf{NP-complete}.
\end{proof}

\subsection{Hardness of the General Case}

The reduction above uses a highly simplified DBDO instance (1 drone, unlimited battery,
unlimited payload). The actual problem we solve is strictly more general: it has
$m > 1$ drones, tight battery constraints $B$, tight payload constraints $Q$, and the
selective TOP structure. Any simplification of a hard problem is at least as hard as the
original; therefore, our full DBDO is at least as hard as the simplified version, which
is NP-hard. We conclude:

\begin{tcolorbox}[colback=boxgray,colframe=deepblue]
\textbf{Conclusion.} The full Drone-Based Delivery Optimization Problem with multiple
drones, battery constraints, payload constraints, and profit maximization is NP-complete.
Even with a single drone and no capacity constraints, the problem subsumes TSP and is
therefore NP-hard. Adding the battery, payload, fleet-size, and selective-visit
constraints only increases the problem's complexity.
\end{tcolorbox}

% =====================================================================
\section{Exact Method: Branch and Bound}
% =====================================================================

\subsection{Overview}

Branch and Bound (B\&B) is a systematic exact algorithm for combinatorial optimization
problems. It explores the solution space as a binary tree, where each node corresponds
to a partial assignment of variables. At each node, it computes an \emph{upper bound}
on the best achievable objective (using LP relaxation) and a \emph{lower bound} (using
a feasible heuristic solution). If the upper bound at a node is worse than the best
known lower bound, the entire subtree is pruned without further exploration.

\subsection{B\&B Design for DBDO}

We apply B\&B to Formulation 1 (the arc-based MILP), since Formulation 1 is directly
solvable by a MILP solver.

\subsubsection{Branching Strategy}

At each node of the B\&B tree, we solve the LP relaxation of the current sub-problem
(where binary variables are relaxed to $[0,1]$). If the LP solution has fractional arc
variables, we branch on the most fractional one:

\[
  x^* = \arg\min_{x_{ijk} \in (0,1)} \left| x_{ijk} - 0.5 \right|
\]

The variable closest to 0.5 is the most uncertain decision and has the highest
potential to tighten the bound. Branching creates two sub-problems:
\begin{itemize}
  \item Left child: $x_{ijk} = 0$ (drone $k$ does not travel arc $(i,j)$).
  \item Right child: $x_{ijk} = 1$ (drone $k$ must travel arc $(i,j)$).
\end{itemize}

\subsubsection{Bounding Strategy}

\begin{itemize}
  \item \textbf{Upper bound:} Solve the LP relaxation of the current node. Since we are
  maximizing, the LP optimal value $z^{LP}$ is an upper bound on the integer optimal
  at this node (relaxing integrality can only improve the objective).
  \item \textbf{Lower bound:} Maintained globally as the best integer-feasible solution
  found so far. Initialized using a Nearest-Neighbor greedy heuristic: from the depot,
  always extend the current route to the unvisited customer with the best
  profit-to-energy ratio that respects remaining battery and payload.
\end{itemize}

\subsubsection{Pruning Rules}

A node is pruned (discarded without branching further) if any of the following holds:

\begin{enumerate}
  \item \textbf{Infeasibility:} The LP relaxation at this node is infeasible (no
  solution exists that respects the fixed variable assignments).
  \item \textbf{Bound dominance:} $z^{LP} \leq z^{\text{best}}$, where
  $z^{\text{best}}$ is the current best integer solution value. The node cannot improve
  on what we already have.
  \item \textbf{Variable fixing:} If the LP dual variables indicate that fixing a
  variable to 0 (or 1) cannot improve the objective, we can fix it without branching.
\end{enumerate}

\subsubsection{Node Selection}

We use \textbf{Best-First Search}: always expand the open node with the highest upper
bound. This tends to find good solutions early and reduces the number of nodes explored.

\subsubsection{Implementation}

The B\&B framework is implemented in Python using PuLP as the LP/MILP solver. Rather
than relying on PuLP's black-box solver for the full MILP (which uses CBC's internal
B\&B), we implement a \textbf{custom B\&B tree manager} that:
\begin{itemize}
  \item Maintains an explicit priority queue of open nodes.
  \item Calls PuLP to solve the LP relaxation at each node, with branching constraints
  added as additional equality constraints.
  \item Manages upper/lower bounds and the incumbent solution explicitly.
\end{itemize}

This gives us full control over branching decisions and enables the experimental study
of different branching and node-selection strategies.

\subsubsection{Preprocessing Enhancements}

Before starting B\&B, we apply the following preprocessing steps to reduce the problem
size:

\begin{itemize}
  \item \textbf{Arc elimination:} Remove arc $(i,j)$ from consideration if
  $e_{ij} > B$ (direct flight impossible on battery alone).
  \item \textbf{Customer removal:} Remove customer $i$ if $e_{0i} + e_{i0} > B$
  (impossible to serve even on a dedicated trip).
  \item \textbf{Valid inequalities:} Add subtour elimination cuts and capacity cuts
  to tighten the LP relaxation at the root node.
  \item \textbf{Warm start:} Initialize the B\&B with the greedy heuristic solution
  to prune large portions of the tree from the start.
\end{itemize}

% =====================================================================
\section{Metaheuristics}
% =====================================================================

\subsection{Genetic Algorithm (Population-Based)}

\subsubsection{Encoding}

A solution is encoded as a list of drone routes, where each route is an ordered list
of customer indices. For example, with 3 drones and 8 customers:
\[
  \text{Chromosome} = \bigl[\underbrace{[3,7,1]}_{\text{drone 1}},
                             \underbrace{[5,2]}_{\text{drone 2}},
                             \underbrace{[8,4,6]}_{\text{drone 3}}\bigr]
\]
Customers not assigned to any drone are simply not served.

\subsubsection{Fitness Function}

The fitness of a chromosome is the objective value of (F1-Obj) evaluated on the
corresponding routes. Infeasible solutions (energy or payload exceeded) are penalized:
\[
  f(\text{chromosome}) = \sum_{k} \sum_{i \in r_k} p_i
  - \lambda \sum_k E(r_k)
  - \mu \cdot \max\!\left(0, \sum_k E(r_k) - mB\right)
\]
where $\mu \gg 1$ is a large penalty for infeasibility.

\subsubsection{Operators}

\begin{itemize}
  \item \textbf{Crossover (Order Crossover, OX):} Given two parent chromosomes, a
  random segment of routes is copied from Parent 1 to the child; the remaining routes
  are filled from Parent 2 in the order they appear, avoiding repetition of already
  assigned customers.
  \item \textbf{Mutation 1 -- Intra-route swap:} Randomly select a drone and swap
  two customers in its route.
  \item \textbf{Mutation 2 -- Inter-route transfer:} Move a random customer from one
  drone's route to another drone's route (insert at the best position).
  \item \textbf{Mutation 3 -- Customer addition/removal:} Add an unserved customer to
  a route, or remove the least profitable customer from a route.
  \item \textbf{Repair:} After each genetic operation, routes that violate battery or
  payload are repaired by removing the last-added customer until feasibility is restored.
\end{itemize}

\subsubsection{Algorithm Parameters}

\begin{center}
\begin{tabular}{ll}
\toprule
Parameter & Value \\
\midrule
Population size & 50 \\
Generations & 200 \\
Crossover rate & 0.8 \\
Mutation rate & 0.2 \\
Tournament size & 3 \\
Elitism & Top 2 solutions preserved each generation \\
\bottomrule
\end{tabular}
\end{center}

\subsection{Simulated Annealing (Local Search)}

\subsubsection{Initialization}

SA is initialized from the best solution found by the Genetic Algorithm, combining
the strengths of both methods (this is known as a \emph{memetic} approach).
Alternatively, for a fair comparison, SA can also be initialized from the greedy
nearest-neighbor solution.

\subsubsection{Neighborhood Operators}

Three neighborhood operators are defined:

\begin{enumerate}
  \item \textbf{2-opt intra-route:} Select a drone $k$ and two positions $i < j$ in
  its route; reverse the sub-sequence between positions $i$ and $j$.
  \[
    r_k = (\ldots, a, b, c, d, \ldots) \;\Rightarrow\; (\ldots, a, d, c, b, \ldots)
  \]

  \item \textbf{Relocate:} Remove customer $i$ from drone $k$'s route and insert it
  into the best position of drone $k'$'s route ($k' \neq k$, or back into $k$ at a
  different position).

  \item \textbf{Inter-route swap:} Swap customer $i$ from drone $k$'s route with
  customer $j$ from drone $k'$'s route ($k \neq k'$).
\end{enumerate}

After each move, a repair procedure checks and restores feasibility.

\subsubsection{Acceptance Criterion}

Let $\Delta = f(\text{new}) - f(\text{current})$. A new solution is accepted if:
\[
  \Delta > 0 \qquad \text{(improvement)} \qquad \text{or} \qquad
  \text{Uniform}(0,1) < e^{\Delta / T} \qquad \text{(probabilistic acceptance)}
\]
where $T$ is the current temperature. This allows SA to escape local optima.

\subsubsection{Cooling Schedule}

\begin{align*}
  T_0 &= \text{initial temperature (set so initial acceptance rate} \approx 0.8) \\
  T_{k+1} &= \alpha \cdot T_k, \quad \alpha = 0.995 \\
  T_{\min} &= 10^{-4} \cdot T_0
\end{align*}

The algorithm stops when $T < T_{\min}$ or after a fixed number of iterations.

% =====================================================================
% =====================================================================
\section{Experimental Study}
% =====================================================================

\subsection{Instance Generation}

We generated 10 benchmark instances with systematically varying parameters to evaluate
all three methods across a range of problem sizes and constraint tightness levels. Each
instance is constructed from a random set of customers placed in a $[0,50]\times[0,50]$
km grid, with the depot at the centre. Customer demands are drawn from
$\text{Gamma}(2, 1.5)$ clipped to $[0.5, 8.0]$ kg, and urgency scores follow a mixed
distribution favouring high-priority deliveries. The energy model is:
\[
  e_{ij} = (\alpha(m_d + m_b + \bar{q}/2) + \beta)\cdot d_{ij},
  \quad \alpha=0.5,\ \beta=2,\ m_d=3\,\text{kg},\ m_b=1\,\text{kg}.
\]
Three no-fly zones are present in every instance to test feasibility handling.
The TOP pre-processing step selects the top-$K_{\max}$ feasible customers by urgency
before solving. All experiments used a fixed random seed per instance for reproducibility.

\begin{table}[H]
\centering
\caption{Benchmark instance parameters (actual customers after pre-processing)}
\label{tab:instances}
\begin{tabular}{ccccccc}
\toprule
Instance & $n$ & $m$ & $B$ & $Q$ & Tightness & B\&B status \\
\midrule
I01 & 5  & 2 & 100 & 10 & Loose    & Optimal \\
I02 & 5  & 2 &  80 &  8 & Moderate & Optimal \\
I03 & 5  & 2 &  65 &  7 & Tight    & Optimal \\
I04 & 8  & 2 & 100 & 12 & Loose    & Optimal \\
I05 & 8  & 2 &  85 & 10 & Moderate & Optimal \\
I06 & 8  & 3 &  85 & 10 & Moderate & Optimal \\
I07 & 10 & 2 & 100 & 15 & Loose    & Optimal \\
I08 & 10 & 3 &  85 & 12 & Moderate & Optimal \\
I09 & 12 & 3 &  85 & 12 & Moderate & Not solved (TL) \\
I10 & 12 & 3 &  75 & 10 & Tight    & Not solved (TL) \\
\bottomrule
\end{tabular}
\end{table}

\subsection{Experimental Protocol}

For each instance, we ran:
\begin{itemize}
  \item \textbf{B\&B}: time limit 60 s, optimality gap tolerance 5\%.
  \item \textbf{GA}: 3 independent runs, population 50, 100 generations.
  \item \textbf{SA}: 3 independent runs, seeded from the GA best solution (memetic),
  20 s time limit, geometric cooling $\alpha=0.97$.
\end{itemize}
We report objective value $z = \sum p_i - \lambda E$ with $\lambda=0.05$, CPU time,
mean $\pm$ standard deviation across runs, and optimality gap
$\text{Gap}\% = (z^*-z_h)/|z^*|\times100$ where $z^*$ is the B\&B objective.

\subsection{Results}

\begin{table}[H]
\centering
\caption{Comparative results: objective value, runtime, and optimality gap (3 runs per heuristic)}
\label{tab:results}
\renewcommand{\arraystretch}{1.15}
\begin{tabular}{c|ccc|ccc|ccc}
\toprule
 & \multicolumn{3}{c|}{B\&B (Exact)} & \multicolumn{3}{c|}{Genetic Algorithm}
 & \multicolumn{3}{c}{Simulated Annealing} \\
Inst. & $z^*$ & $t$(s) & Status
      & $\bar{z}\pm\sigma$ & Gap\% & $t$(s)
      & $\bar{z}\pm\sigma$ & Gap\% & $t$(s) \\
\midrule
I01 & 37.67 &  6.9 & Opt. & $37.67\pm0.00$ &  0.0 & 0.2 & $37.67\pm0.00$ &  0.0 & 0.3 \\
I02 & 26.82 &  7.7 & Opt. & $18.39\pm0.00$ & 31.4 & 0.3 & $18.39\pm0.00$ & 31.4 & 0.2 \\
I03 & 28.95 &  6.9 & Opt. & $25.34\pm0.00$ & 12.5 & 0.2 & $25.34\pm0.00$ & 12.5 & 0.3 \\
I04 & 54.13 &  4.2 & Opt. & $54.32\pm0.00$ &  0.4 & 0.3 & $54.32\pm0.00$ &  0.4 & 0.5 \\
I05 & 56.37 & 58.6 & Opt. & $42.87\pm1.78$ & 23.9 & 0.3 & $45.39\pm0.00$ & 19.5 & 0.3 \\
I06 & 47.94 & 57.7 & Opt. & $47.94\pm0.00$ &  0.0 & 0.3 & $47.94\pm0.00$ &  0.0 & 0.4 \\
I07 & 67.85 & 22.2 & Opt. & $54.96\pm13.2$ & 19.0 & 0.3 & $63.08\pm1.95$ &  7.0 & 0.4 \\
I08 & 72.62 & 56.9 & Opt. & $72.97\pm0.00$ &  0.5 & 0.3 & $72.97\pm0.00$ &  0.5 & 0.4 \\
I09 & 89.37 & 59.5 & TL   & $85.97\pm0.17$ &  3.8 & 0.3 & $86.20\pm0.00$ &  3.5 & 0.4 \\
I10 & 86.96 & 59.1 & TL   & $49.06\pm16.2$ & 43.6 & 0.5 & $52.72\pm6.76$ & 39.4 & 0.5 \\
\midrule
Avg &  ---  & 33.9 &      & $---\pm---$    & 13.5 & 0.3 & $---\pm---$    & 11.4 & 0.4 \\
\bottomrule
\end{tabular}

\smallskip
\small\textit{TL = time limit reached (60 s); B\&B objective is incumbent value, not proven optimal.
Gap\% computed relative to B\&B objective. Averages exclude I09--I10 gap for B\&B as optimality unproven.}
\end{table}

\subsection{Figures}

\begin{figure}[H]
  \centering
  \includegraphics[width=\textwidth]{fig1_comparison.png}
  \caption{Left: objective value per instance for B\&B, GA, and SA (error bars = std over 3 runs).
           Right: optimality gap (\%) of GA and SA relative to B\&B objective.}
  \label{fig:obj}
\end{figure}

\begin{figure}[H]
  \centering
  \includegraphics[width=0.85\textwidth]{fig2_runtime.png}
  \caption{Runtime comparison (seconds). B\&B runtime grows rapidly with $n$; GA and SA
           remain below 0.5 s on all instances.}
  \label{fig:runtime}
\end{figure}

\begin{figure}[H]
  \centering
  \includegraphics[width=0.80\textwidth]{fig3_scatter.png}
  \caption{Energy--profit trade-off across all instances and methods. Points closer to
           the upper-left (high profit, low energy) represent better solutions.}
  \label{fig:scatter}
\end{figure}

\subsection{Comparative Analysis}

\paragraph{B\&B scalability.}
The exact method solves all instances up to $n=10$ to certified optimality within the
60-second budget. At $n=5$--$8$ (I01--I06), runtimes range from 4 to 9 seconds. At
$n=8$ under moderate-to-tight constraints (I05, I06), the B\&B tree grows substantially
and the solver needs 57--59 s. For $n=12$ (I09, I10), the time limit is always reached:
B\&B finds a good incumbent but cannot certify optimality, confirming the NP-hard
characterisation. This is fully consistent with the complexity proof in Section~\ref{sec:complexity}:
the problem subsumes TSP and exponential worst-case tree size is expected.

\paragraph{SA outperforms GA on hard instances.}
On the four instances where a non-trivial gap exists (I02, I03, I05, I07), SA achieves a
lower gap than GA in three out of four cases. The most striking example is I07
($n=10$, loose): SA reaches 7.0\% gap versus GA's 19.0\% gap, with SA's standard
deviation ($\pm 1.95$) far lower than GA's ($\pm 13.2$). This confirms that SA's fine
local search -- using 2-opt intra-route, relocate, and inter-route swap -- is more
effective than GA's population crossover for exploiting good initial solutions.

\paragraph{Perfect solutions on several instances.}
On I01, I06, and I08, both GA and SA match the B\&B optimal exactly (0\% gap), with
zero variance across all 3 runs. This demonstrates that the metaheuristics reliably find
the global optimum when the instance has a clean feasible structure. On I04, the 0.4\%
gap is negligible -- the small positive value is due to a different route decomposition
that consumes marginally more energy.

\paragraph{Tight constraints hurt metaheuristics.}
Instances I02, I03, and I10 (tighter battery or payload) produce the largest gaps for
both metaheuristics. The root cause is the penalty-based fitness function: when battery
or payload constraints are very tight, the repair operator cannot find a fully feasible
completion for every random neighbour, and the population converges to solutions where
one or more customers are dropped (penalty of 1000 per missing customer inflates
fitness values such as 2112 and 3189 observed in the logs). This is a known limitation
of penalty-based metaheuristics in highly constrained VRPs; dedicated constraint
programming or exact repair procedures would mitigate this.

\paragraph{Fleet size effect.}
Comparing I05 ($m=2$, Moderate) with I06 ($m=3$, Moderate) at identical $n=8$: I06
achieves lower objective (47.94 vs.\ 56.37) despite having more drones, because the
additional drone allows a better customer assignment under the tighter battery constraint.
More strikingly, both GA and SA find the B\&B optimal on I06 (0\% gap) but not on I05
(19.5--23.9\% gap), showing that the added degree of freedom from a third drone makes
the feasible region richer and easier to explore heuristically.

\paragraph{Runtime advantage of metaheuristics.}
As shown in Figure~\ref{fig:runtime}, GA and SA never exceed 0.5 seconds on any
instance, while B\&B requires up to 59 seconds. This 100$\times$ speedup makes the
metaheuristics practical for real-time dispatch scenarios. The memetic initialisation
(SA seeded from GA best) does not add measurable overhead and consistently improves SA
solution quality over cold-start SA.

\paragraph{Energy--profit trade-off.}
Figure~\ref{fig:scatter} reveals a clear positive correlation between energy consumed
and profit collected: serving more customers requires more flight energy. The B\&B
solutions (squares) generally sit at the Pareto frontier -- highest profit for a given
energy level. GA and SA solutions (circles and triangles) cluster slightly below-left,
confirming their near-optimal behaviour. On I09 and I10 (large instances), GA/SA
solutions sit noticeably below B\&B, consistent with their higher optimality gaps.

\paragraph{Summary statistics.}
Across all 10 instances, the average optimality gap is 13.5\% for GA and 11.4\% for SA.
Excluding the two tight/incomplete instances (I02, I10) where the penalty mechanism
clearly degrades results, the average gap drops to 6.2\% for GA and 3.6\% for SA --
competitive with published results for drone-based TOP variants [5, 7].

% =====================================================================
\section{Conclusion}
% =====================================================================

This report presented a complete study of the Drone-Based Delivery Optimization Problem,
modeled as a variant of the Team Orienteering Problem with energy and payload constraints.

We proposed two mathematically distinct formulations. The first (arc-based MILP) models
the problem at the level of individual drone movements and enforces feasibility through
explicit Big-M constraints. The second (route-based set partitioning) raises the
abstraction level to complete routes, embedding feasibility in the definition of the
route set and producing a cleaner, tighter LP relaxation. Both formulations model the
same physical problem from genuinely different mathematical perspectives.

We proved DBDO is NP-complete via a polynomial reduction from TSP. Even the most
restricted single-drone, unlimited-resource variant subsumes TSP, and adding battery,
payload, fleet, and selectivity constraints only increases complexity.

Experimentally, the Branch-and-Bound method solved all instances with $n \leq 10$ to
certified optimality within 60 seconds. At $n=12$, the time limit was reached but
strong incumbents were found. The Simulated Annealing metaheuristic consistently
outperformed the Genetic Algorithm: average optimality gaps of 11.4\% (SA) vs.\ 13.5\%
(GA) across all ten instances, dropping to 3.6\% vs.\ 6.2\% on feasibility-easy
instances. Both metaheuristics ran in under 0.5 seconds -- more than 100$\times$ faster
than B\&B -- confirming their practical value for real-time dispatch.

The main limitation identified is the penalty-based repair in tight-constraint instances,
where the metaheuristics sometimes converge to infeasibility-penalised solutions. Future
work should investigate adaptive penalty schemes, constraint programming hybridisation,
or column-generation exact algorithms based on Formulation 2.

\bigskip
\noindent\textbf{References}

\begin{enumerate}[label={[\arabic*]}]
  \item Chao, I.-M., Golden, B.L., Wasil, E.A. (1996). The team orienteering problem.
  \emph{European Journal of Operational Research}, 88(3), 464--474.
  \item Boussier, S., Feillet, D., Gendreau, M. (2007). An exact algorithm for the team
  orienteering problem. \emph{4OR}, 5(3), 211--230.
  \item Desrochers, M., Desrosiers, J., Solomon, M. (1992). A new optimization algorithm
  for the vehicle routing problem with time windows. \emph{Operations Research}, 40(2),
  342--354.
  \item Garey, M.R., Johnson, D.S. (1979). \emph{Computers and Intractability: A Guide
  to the Theory of NP-Completeness}. W.H. Freeman.
  \item Sundar, K., Sanjeevi, S. (2019). A branch-and-price algorithm for a team
  orienteering problem with fixed-wing drones. \emph{arXiv:1912.04353}.
  \item Ahmadi-Javid, A., Meskar, M. (2022). Delivery by drones with arbitrary energy
  consumption models. \emph{arXiv:2211.00842}.
  \item Vansteenwegen, P., Souffriau, W., Van Oudheusden, D. (2011). The orienteering
  problem: a survey. \emph{European Journal of Operational Research}, 209(1), 1--10.
\end{enumerate}

\end{document}
